% ============================================================
%  LAMPIRAN 8 – PROFIL DAN KUALIFIKASI PAKAR
%  Compile standalone:  pdflatex lampiran-8.tex
%  Compile as part of thesis: pdflatex main.tex
% ============================================================
\documentclass[../../thesis]{subfiles}
\usepackage{hyphenat}
\usepackage{iftex}

% ── Encoding, Language & Fonts ───────────────────────────────
\ifXeTeX
    \usepackage{polyglossia}
    \setmainlanguage{bahasai}
    \usepackage{fontspec}
    \setmainfont{TeX Gyre Termes}
    \setsansfont{TeX Gyre Heros}
    \setmonofont[Scale=0.88]{TeX Gyre Cursor}
\else
    \usepackage[utf8]{inputenc}
    \usepackage[T1]{fontenc}
    \usepackage[indonesian]{babel}
    \usepackage{mathptmx}
    \usepackage{helvet}
    \usepackage{courier}
\fi

% ── Page geometry (A4, UI thesis margins) ────────────────────
\usepackage[
  a4paper,
  headheight = 0pt,
  top    = 3cm,
  bottom = 3cm,
  left   = 3cm,
  right  = 3cm
]{geometry}

% ── Line spacing & paragraph ─────────────────────────────────
\usepackage{setspace}
\onehalfspacing
\setlength{\parindent}{1.27cm}
\setlength{\parskip}{0pt}

% ── Microtype & justification ────────────────────────────────
\usepackage{microtype}
\sloppy

% ── Headings ─────────────────────────────────────────────────
\usepackage{titlesec}

\titleformat{\chapter}[block]
  {\centering\rmfamily\bfseries\large}{}
  {0pt}{\MakeUppercase}
\titlespacing*{\chapter}{0pt}{0pt}{1.5em}

\titleformat{\section}
  {\rmfamily\bfseries\normalsize}{\thesection}{1em}{}
\titlespacing*{\section}{0pt}{1.2em}{0.6em}

\titleformat{\subsection}
  {\rmfamily\bfseries\normalsize}{\thesubsection}{1em}{}
\titlespacing*{\subsection}{0pt}{1em}{0.5em}

% ── Tables ───────────────────────────────────────────────────
\usepackage{longtable, booktabs, array, multirow}
\usepackage{calc}
\usepackage{tabularx}

% ── Hyperlinks ───────────────────────────────────────────────
\usepackage[hidelinks]{hyperref}

% Bibliography setup is inherited from main.tex (biblatex + references.bib)
% when this file is compiled standalone via the subfiles class.

% ============================================================
\begin{document}

% When standalone: switch to appendix mode and set the counter so the
% appendix heading renders correctly. When compiled inside main.tex,
% main.tex issues \appendix and advances the counter before \subfile.
\ifSubfilesClassLoaded{}{\appendix\setcounter{chapter}{1}}

\lampiran{Profil dan Kualifikasi Pakar}
\label{lamp:profil-pakar}

Lampiran ini merinci latar belakang keilmuan panel pakar yang memvalidasi
keluaran sistem (Subbab~\ref{subsec:subjek-evaluasi}). Identitas pakar
disajikan secara anonim dalam bentuk kualifikasi keilmuan---bidang studi,
institusi, dan tahun---tanpa menyebut nama pribadi. Penetapan pakar mengikuti
kriteria \textit{purposive sampling} pada Bab~3, yakni mahasiswa aktif atau
lulusan dengan latar belakang ilmu murni atau kependidikan yang linier dengan
mata pelajaran yang dinilai.

\begin{longtable}{@{}
  >{\raggedright\arraybackslash}p{0.16\linewidth}
  >{\raggedright\arraybackslash}p{0.16\linewidth}
  >{\raggedright\arraybackslash}p{0.60\linewidth}@{}}
\caption{Profil keilmuan enam pakar validasi (dua per mata pelajaran). Label
  ``Expert~1''/``Expert~2'' bersifat lokal per mata pelajaran, konsisten dengan
  Tabel~\ref{tab:validasi-mapel}.}
\label{tab:profil-pakar}\\
\toprule
\textbf{Mata Pelajaran} & \textbf{Pakar} & \textbf{Latar belakang keilmuan} \\
\midrule
\endfirsthead
\toprule
\textbf{Mata Pelajaran} & \textbf{Pakar} & \textbf{Latar belakang keilmuan} \\
\midrule
\endhead
\bottomrule
\endlastfoot
\multirow{2}{*}{Biologi}
  & Expert 1 & Pendidikan Biologi, Universitas Pendidikan Indonesia (2022) \\
  & Expert 2 & Biologi, Universitas Indonesia (2022) \\
\midrule
\multirow{2}{*}{Fisika}
  & Expert 1 & Pendidikan Fisika, Universitas Pendidikan Indonesia (2023) \\
  & Expert 2 & Fisika, Universitas Gadjah Mada (2022) \\
\midrule
\multirow{2}{*}{Kimia}
  & Expert 1 & Pendidikan Kimia, Universitas Pendidikan Indonesia (2022) \\
  & Expert 2 & Kimia, Institut Teknologi Bandung (2022) \\
\end{longtable}

Komposisi dua pakar per mata pelajaran sengaja memadukan latar \emph{kependidikan}
(Expert~1) dan \emph{ilmu murni} (Expert~2) agar penilaian mencakup ketepatan
pedagogis sekaligus akurasi keilmuan. Susunan ini juga memungkinkan perhitungan
kesepakatan antarpakar (Gwet's AC1) pada validasi Fase~1.

Pada validasi lintas-disiplin (Fase~2), penilaian dilakukan oleh panel diskusi
yang terdiri atas \textbf{tiga pakar}, yakni satu wakil per mata pelajaran
(Expert~2 dari Biologi, Fisika, dan Kimia). Ketiga pakar lintas-disiplin ini
menilai relasi lintas-buku secara bersama untuk mencapai satu jawaban konsensus
per relasi, sehingga, berbeda dari Fase~1 yang melibatkan dua pakar independen
per \textit{triple}, kesepakatan antarpakar tidak dihitung pada fase ini.

\end{document}
